% mage-auto-heading: modeling-is-a-classical-engineering-move
% mage-auto-heading: a-model-is-a-purposeful-reduction
% mage-auto-heading: tolerance-and-degrees-of-freedom
% mage-auto-heading: not-every-unmodeled-choice-is-free
% mage-auto-heading: constructing-and-revising-a-model
% mage-auto-heading: implementation-can-induce-a-model
% mage-auto-heading: judging-what-an-agent-discovers
% mage-auto-heading: structure-and-semantics-make-representations-machine-operable
% mage-auto-heading: representations-support-accumulating-capabilities
% mage-auto-heading: the-engineering-modeling-repertoire
% Authoritative LaTeX manuscript; Markdown is a generated compatibility projection.

% <!-- part-title: Modeling -->

% <!-- chapter-title: From Context to Engineering Models -->

\chapter{From Context to Engineering Models}

% <!-- section-terms: modeling-principle, reasoning-horizon, context-window -->

% <!-- point: a-reasoner-arrives-knowing-nothing-about-your-system | A foundation model knows the languages and patterns; it does not know your system — those facts live in the engineered environment. | terms: reasoning-horizon, context-window, modeling-principle -->

A foundation model brings knowledge of programming languages, frameworks, databases, distributed systems, and common engineering practices. Fresh out of the box, it does not know the particular system in front of it: its boundaries, history, conventions, ownership, or current state. That knowledge is distributed across source code, configuration, documentation, registries, version history, and the people who maintain the system. An agent can reconstruct what a task needs from those sources, but doing so takes work. If the result is not preserved, the next agent may have to repeat that work for a closely related task.

% <!-- box-family: inset -->

% <!-- inset-domain: AGENT HARNESS -->

% <!-- box-weight: aside -->

% mage-box-prefix: Inset —

\begin{bookbox}{inset}{Skills and Project Instructions}
\label{inset-skills-and-project-instructions}
MAGE\textquotesingle s emphasis on externalizing engineering knowledge appears in today\textquotesingle s agent harnesses as project instructions and skills. Project instructions keep small amounts of cross-cutting, repository-specific guidance persistently available; skills package larger procedures and resources for a recurring class of work and load them when that work arises. Both reduce what a person must supply or an agent must reconstruct in each new reasoning episode.

The progression continues past prose. Retrieval can select knowledge from a larger external store, and models can represent it in structures that support queries, transformations, validation, or deterministic execution. As the representation becomes more structured, the environment can do more with the knowledge than simply show it to the agent.

Skills and instruction files are contemporary carriers. The engineering requirement survives changes in harness design: consequential knowledge that would otherwise be tacit or repeatedly reconstructed should be externalized in a form suited to how it will be used.
\end{bookbox}

% <!-- point: the-first-purpose-of-a-model-is-better-context | The first modeling move is simple: externalize engineering knowledge that should not have to be rediscovered. | terms: modeling-principle, reasoning-horizon, context-window -->

The first \textbf{Modeling move} is therefore to make reusable engineering knowledge explicit. Preserve the entities and relationships that recurring engineering questions depend on so that a task can retrieve the connected slice it needs rather than reconstructing that structure from raw artifacts.

% <!-- point: natural-language-is-itself-a-representation | Natural language is a representation, and often an appropriate one; consequential knowledge kept only in prose must be articulated correctly and then recovered correctly by each later reasoner, and a structured model can preserve distinctions that would otherwise be restated or reconstructed. | terms: modeling-principle, structured -->

The form of the context matters as well as its size. Natural language is often an appropriate representation: it is inexpensive to write, expressive, and sufficient for many engineering questions. But when we rely on prose, two things must go right. The engineer must state the relevant distinctions clearly, and each later reader---human or agent---must interpret them correctly. A structured model can record some of those distinctions explicitly, so that later work can use them without reconstructing them from the prose each time. This does not make the model correct by construction. It changes which distinctions the engineer must make explicit and which the reader must infer. \magevalue{chapter:theory} examines these two sources of error and their costs.

% <!-- point: models-are-also-for-the-engineer | Models help engineers make decisions and detect omissions as well as helping agents interpret engineering knowledge. | terms: modeling-principle -->

Models also help the engineer decide what to require. Choosing entities, relationships, and permitted values can expose omissions that are easy to overlook in prose. A representation may reveal that two cases need different treatment, or that a decision has not yet been made. That benefit arises while the engineer constructs and examines the model, before any agent uses it. Some knowledge may still be best expressed in prose, and some decisions may still require human judgment.

% <!-- point: the-modeling-principle-in-its-weakest-form | The Modeling Principle in its weakest form: externalize engineering knowledge into explicit, structured models to reason through. | terms: modeling-principle, reasoning-horizon, context-window -->

% <!-- index-def: modeling-principle -->

\textbf{The Modeling Principle:} externalize engineering knowledge and intent into explicit, structured models that both engineers and agents can reason through. In practice, give recurring engineering questions purposeful representations. A useful representation may bring a question within the reasoner's capacity, or make an already manageable question cheaper to answer.

Software is unusually changeable. As a system evolves, engineers need to track what should change and which characteristics must be preserved. Requirements engineering helps engineers state what stakeholders want and distinguish those goals from assumptions about the environment\citep{zave1997darkcorners}\citep{nuseibeh2000roadmap}\citep{vanlamsweerde2009re}. Requirements may initially be qualitative or incomplete, and can depend on assumptions about the environment. Making requirements explicit lets engineers examine and refine them, and assess whether a later change preserves the properties that matter.

Engineers also develop specifications that state what an implementation must do. An interface contract may define valid inputs and promised results; a state machine may define permitted transitions. These representations let engineers reason about whether a design satisfies its requirements while leaving other implementation choices open. MAGE draws on these practices. This chapter examines how to choose representations that help engineers and agents reason about a system. \magevalue{chapter:alignment} then examines which obligations the engineered environment should enforce.

\section{Modeling Is a Classical Engineering Move}\label{heading:modeling-is-a-classical-engineering-move}

% <!-- point: modeling-is-a-classical-engineering-move-not-a-mage-invention | Modeling is not new: engineering has long used representations that suppress irrelevant detail so a property can be reasoned about directly. | terms: modeling-principle -->

Modeling is not new. Engineering disciplines have long used representations that suppress irrelevant detail so that particular properties of a system can be reasoned about directly, and software engineering inherited and developed the same practice --- architectural descriptions, schemas, state machines, interface contracts, formal models. Model-based systems engineering makes this commitment particularly explicit: engineering knowledge is organized around models that support analysis, communication, design, and assurance across the system lifecycle\citep{incose2023handbook}\citep{isoiec42010}.

% <!-- point: mage-reuses-established-modeling-traditions | These representations share one engineering move — choose a form in which the property of interest becomes easier to express, inspect, or analyze than in the implementation; MAGE proposes no new modeling language, and this chapter illustrates the technique through selected established models. | terms: modeling-principle -->

These representations differ in purpose and rigor, but they share the same engineering move: choose a form in which the property of interest becomes easier to express, inspect, or analyze than it is in the implementation. UML standardized notations for many structural and behavioral views,\citep{booch-uml} while requirements engineering developed representations for goals, assumptions, and specifications\citep{vanlamsweerde2009re}. MAGE uses these established representations, and this chapter illustrates the technique through selected models.

% <!-- point: mage-adopts-the-classical-move-under-changed-economics | Established models preserve selected design knowledge outside implementation so agents can reuse it; each representation must justify its upkeep. | terms: modeling-principle, alignment-principle -->

Much ordinary software development has relied on a little whiteboarding followed by implementation, with experienced programmers carrying the remaining design knowledge in their heads. That can be an economical arrangement when the programmers themselves perform the implementation. It becomes less attractive when many agents need to recover the same knowledge independently. MAGE uses established modeling techniques to preserve selected engineering decisions outside the implementation, while agents help construct, maintain, and query those representations. The engineer still has to decide which models justify their upkeep.

Models also participate in architectural reasoning. They expose relationships and make selected consequences available for analysis; engineers decide the decomposition, boundaries, interactions, allocations, and tradeoffs the system should embody. A model may record an architectural decision already made or reveal that the architecture should change.

% <!-- box-family: inset -->

% <!-- inset-domain: SOFTWARE ARCHITECTURE -->

% mage-box-prefix: Inset —

\begin{bookbox}{inset}{Models, views, and software architecture}
\label{inset-models-views}
MAGE follows Fairbanks\textquotesingle s definition of software architecture as the set of structures needed to reason about a system, comprising software elements, relations among them, and properties of both\citep{fairbanks2010}. Engineers therefore reason about one system through multiple views: a component view may expose decomposition, a behavioral view legal transitions, a deployment view allocation, and a quantitative view a resource bound. No one view need contain everything another does.

Architectural reasoning may therefore cross several model classes. An architectural decision may depend on structural, behavioral, ownership, decision, measurement, or provenance relationships together. The useful views follow from the engineering questions that must be answered.

Software makes the relationship between model and implementation unusually fluid because both are made of information. A schema can generate code; a state machine can drive behavior; a policy model can generate configuration; a dependency graph can become an execution plan. Such a representation can participate directly in the system while remaining selective about what it represents.

When structured views are cheap to derive, maintain, connect, query, analyze, and sometimes execute, more of the knowledge required for architectural reasoning can remain explicit rather than being repeatedly reconstructed from implementation.
\end{bookbox}

% <!-- point: a-useful-model-need-not-be-elaborate | A useful model may be three states and two transitions, a list of permitted edges, or a schema joined on identifiers; the relevant question is what representation makes a recurring engineering question easier to answer. | terms: modeling-principle -->

A useful model need not be elaborate. Sometimes it is three states and two legal transitions; sometimes a list of permitted service edges; sometimes a schema joined on stable identifiers. The relevant question is what representation makes a recurring engineering question easier to answer.

\section{A Model Is a Purposeful Reduction}\label{heading:a-model-is-a-purposeful-reduction}

% <!-- section-terms: modeling-principle, model-as-map, scope-of-modeling -->

% <!-- point: a-model-is-a-reduction-not-an-approximation | A model is a purposeful reduction; its value is what it leaves out. The engineering question decides what to keep. | terms: model-as-map, scope-of-modeling, modeling-principle -->

% <!-- index-def: model -->

A \textbf{model} is a purposeful reduction of a system. It preserves the information one engineering question needs and suppresses the detail that question does not. Its value is not completeness but selective omission. A reduction an engineer or agent can reason through is more useful than a faithful copy too large to hold in view.

Purposeful reduction can change the difficulty and character of reasoning. A property that requires semantic reconstruction over thousands of implementation details may become a relation over a small graph, a transition over a finite state space, or a bound over a measured quantity. We gain analytical power without having to construct or examine the complete implementation. Engineers still have to establish whether the model’s conclusions are trustworthy and decide what to do with them.

% <!-- point: a-reduction-must-preserve-the-distinctions-its-question-requires | A model can record accurate facts yet omit a distinction needed for a decision: remaining prepaid credits cannot distinguish an exhausted allowance from unlimited service entitlement. | terms: model-as-map, modeling-principle -->

Reduction is useful only if it preserves the distinctions the engineering question requires. When it doesn't, cases that the system must treat differently can become indistinguishable in the model. Suppose a service offers prepaid credits and unlimited subscriptions. Before accepting a request, the system must determine whether the customer is entitled to use the service. A model that records only remaining prepaid credits gives the same balance---zero---to a prepaid customer who has spent every credit and an unlimited subscriber who never needed to buy any. A decision rule that admits only customers with a positive balance would then wrongly refuse the unlimited subscriber. The balances are accurate, but the model has lost the distinction between an exhausted allowance and unlimited use.

This is the connection Modeling shares with representation learning: what a reasoner can infer efficiently depends partly on the form in which the problem is presented. MAGE begins from representations engineering has already developed for consequential system questions rather than asking a learned reasoner to rediscover those abstractions on every task.

% <!-- point: reduction-preserves-degrees-of-freedom | Obligations rarely fix one implementation; a tolerance bounds permitted variation, a degree of freedom is a choice left open. | terms: model-as-map, scope-of-modeling, degrees-of-freedom, tolerance -->

\subsection{Tolerance and degrees of freedom}\label{heading:tolerance-and-degrees-of-freedom}

Leaving detail out of a model can also be understood as preserving useful choices during implementation. Engineering obligations rarely determine a single acceptable result. They constrain some variation and leave other choices open. This book uses two engineering terms for that distinction:

\begin{itemize}
\tightlist
\item
  A \textbf{tolerance} bounds the variation an obligation permits.
\item
  A \textbf{degree of freedom} is a choice left open because the governing obligations do not distinguish among its acceptable alternatives.
\end{itemize}

A tolerance can restrict an implementation choice without determining it completely. Purposeful reduction therefore does two things at once: it represents the distinctions the engineering question requires while declining to settle choices that do not matter to that question. In the MAGE methodology, engineers leave those degrees of freedom open unless an engineering obligation gives them a reason to constrain the choice.

For example, a model can record that a list of up to 1,000 names must be returned in alphabetical order within 200 milliseconds. The timing requirement sets a tolerance: response times may vary as long as they stay within that bound. If several sorting algorithms meet the ordering and timing requirements, choosing among them is a degree of freedom left to the engineer or agent. \autoref{realization-space} shows how the ordering requirement and response-time tolerance bound the acceptable implementations while leaving choices among them.

\begin{figure}[htbp]
\centering
\includegraphics[width=\linewidth]{book/latex/assets/realization-space-2a3916a20a.pdf}
\caption{\emph{Requirements bound the acceptable implementation space.} Each dot represents a possible implementation for lists of up to 1,000 names. Alphabetical ordering defines the left region, and the 200-millisecond response-time tolerance defines the right. Their intersection contains multiple acceptable implementations, leaving degrees of freedom such as the choice of algorithm.}
\label{realization-space}
\end{figure}

% <!-- point: freedom-defers-commitment | Preserving a degree of freedom defers a decision to a reasoner better positioned to make it: information may arrive later, and the implementing agent can acquire evidence before deciding. | terms: degrees-of-freedom, scope-of-modeling -->

Why might it be helpful to preserve degrees of freedom? Let\textquotesingle s look at two different reasons. First, information may arrive later. The target environment, surrounding implementation, workload, available components, or consequences of earlier choices may not be known when the model is constructed. A decision made during implementation can use information that an earlier decision could not. Second, the implementing agent can actively acquire evidence before deciding. It can construct alternatives, run pilots, benchmark them, measure their consequences, and tune the implementation to its environment. Human engineers can of course do the same, as can evolutionary algorithms, autotuners, search procedures, and other optimization mechanisms. What changes with commodity intelligence is how broadly and cheaply this kind of exploration can be applied during ordinary implementation. Closing a degree of freedom early can therefore replace a later, evidence-informed decision with an earlier guess. Engineering should constrain the choices whose consequences matter now while preserving choices that can be made better during implementation.

% <!-- point: parsimony-is-the-complementary-bottom-up-test | Parsimony reasons from the proposed model inward: each distinction it contains must be required by the engineering question or removed; the aim is not the smallest model but the smallest model adequate to its purpose. | terms: degrees-of-freedom -->

Parsimony---keeping only what is needed---gives us a second way to examine a proposed model. We began with the engineering question and asked what the model must preserve and what it may leave unspecified. Now consider each distinction already in the model: which question or obligation requires it? If it serves no such purpose, removing it can reduce both the engineer's work and the work of keeping the model current. The aim is the smallest model adequate to its purpose.

% <!-- point: reduction-fails-in-either-direction | Reduction fails in either direction: too little reproduces complexity and may needlessly constrain implementation, while too much erases a consequential distinction. | terms: degrees-of-freedom, modeling-principle -->

A model can fail on either side of that aim:

\begin{itemize}
\tightlist
\item
  \textbf{Too little reduction} retains unnecessary distinctions. They consume attention, add model maintenance work, and may constrain implementation without improving the answer.
\item
  \textbf{Too much reduction} removes a distinction needed to answer the question. Recall the earlier subscription example: the credit balance could not distinguish an exhausted prepaid allowance from unlimited service.
\end{itemize}

% <!-- point: mage-prefers-adequate-reductions-not-more-models | MAGE favors models adequate to their purpose; both additional detail and greater semantic precision must justify their construction and maintenance costs. | terms: modeling-principle, degrees-of-freedom, semantic-commitment -->

A parsimonious model preserves the consequential distinctions and provides enough explanation for its intended users to work with them. Adding detail is therefore not always an improvement. Sometimes improving a model means deleting structure that no longer helps answer its question. The same test applies to precision of meaning, which we examine below: the benefit must justify the effort of defining and maintaining it.

% <!-- point: models-are-provisional-engineering-judgments | Whether a distinction is needed remains an engineering judgment that can change with evidence; an omitted distinction may prove consequential, or a represented one unnecessary. | terms: degrees-of-freedom, modeling-principle -->

These judgments can change with experience. An omitted distinction may turn out to matter, or the system may change so that a distinction once needed is no longer useful. A model records our present understanding of what matters, and we revise it as that understanding improves.

% <!-- point: software-tolerance-is-often-qualitative | In software a tolerance can be numeric but is often qualitative, because code owes more to discrete mathematics than to continuous media. | terms: tolerance, degrees-of-freedom -->

In software engineering, a tolerance can be numeric but is often qualitative. A latency budget may set a numerical bound: a response must complete within 200 ms. Coupling admits an intermediate form: an architecture may tolerate dependencies up to a measured threshold while also distinguishing kinds of coupling that are acceptable or prohibited. Other architectural tolerances are primarily qualitative: one component may depend on another through a declared interface but not reach directly into its implementation. The difference partly follows from the medium. Physical engineering often concerns continuous quantities such as length, voltage, pressure, or temperature, so tolerances naturally take numerical form. Software structures and behaviors are largely discrete---types, sets, relations, graphs, states, and transitions---so tolerances often define sets of acceptable implementations instead.

% <!-- point: a-measurement-model-can-carry-a-margin-not-just-a-target | A quantitative model can distinguish a target from a required limit, allowing engineers to reason about the remaining margin as well as whether the requirement is satisfied. | terms: tolerance -->

A quantitative model can also distinguish a target from a limit. A service may aim for latency below its required maximum, leaving a margin for variation. A capacity model may likewise distinguish expected demand from the capacity available to meet it. Keeping those quantities separate lets engineers reason about the remaining margin as well as whether a requirement is currently satisfied.

% <!-- point: obligations-bound-a-region-without-choosing-a-point | Obligations bound the space of acceptable implementations without selecting a point; commodity intelligence makes exercising that freedom cheap, so Modeling names the consequential distinctions and Alignment enforces selected boundaries. | terms: degrees-of-freedom, tolerance -->

Software also admits considerable variation in how a requirement is implemented. Programs are constrained less by physical materials than by the information they must accept, produce, and preserve. Many implementations may therefore satisfy the same obligations\citep{brooks1987}. Commodity intelligence makes it cheaper to explore those alternatives. An agent can construct, test, and revise implementations within the acceptable region, while the engineer determines which properties bound that region.

% <!-- point: capability-moves-the-specification-boundary-not-the-obligation | The engineering problem sets the obligations regardless of the implementer; a more capable human or agent may need less scaffolding but still needs the relevant information in the model or from other evidence. -->

How much must be written down also depends on who will implement it. An engineer or agent may need examples, a decomposition, or intermediate constraints to interpret and satisfy an obligation reliably. These provide scaffolding in addition to the obligation itself. \autoref{specified-region} distinguishes the two.

Let me emphasize: the system's obligations are set by the engineering problem, not by the implementer. Greater capability does not make the requirements optional. However, a more capable actor, whether human or agent, may need less scaffolding to meet them, allowing us to leave more implementation choices open. The actor still needs the information required to meet those obligations, either in the model or from other evidence.

\begin{figure}[htbp]
\centering
\includegraphics[width=\linewidth]{book/latex/assets/specified-region-bfae6b28e3.pdf}
\caption{\emph{Obligation, scaffolding, and degrees of freedom.} The engineering obligation establishes the minimum consequential content that must be preserved. A particular agent may require additional scaffolding to interpret and satisfy that obligation reliably. Greater agent capability can reduce this scaffolding and leave more implementation choices open, but it does not change the obligation itself.}
\label{specified-region}
\end{figure}

% <!-- point: tacit-obligation-vs-genuinely-free | A tacit obligation can look free; the "not that" test reveals a boundary that was never written down. | terms: degrees-of-freedom, tolerance, modeling-principle -->

\subsection{Not every unmodeled choice is free}\label{heading:not-every-unmodeled-choice-is-free}

This also distinguishes tacit obligations from genuinely free choices. An unmodeled choice can look open even when engineers would reject some implementations of it. Then the choice was never wholly free. Human judgment was already imposing a boundary. What was missing was an explicit representation of the obligation or its tolerance. A genuinely free choice is different. Once the governing obligations are accounted for, engineering has no reason to distinguish among the remaining alternatives. A useful diagnostic follows: if an implementation would cause the reviewer to say, ``\textbf{\emph{NOT THAT!}}'', then some consequential boundary remains, whether or not anyone has written it down.

% <!-- point: unknown-tacit-explicit-free-are-states | Whether a choice is tacitly bounded or genuinely free is a present state, not a fixed one; it changes as engineers learn from the system. | terms: degrees-of-freedom, modeling-principle -->

Whether a choice is tacitly bounded or genuinely free can change as engineers learn from the system. A region initially treated as free may expose a consequential distinction; investigation may also show that a suspected distinction does not matter.

% <!-- point: modeling-can-precede-complete-specification | Modeling can begin before specification is complete; a reduction can make variation tractable so distinctions can be discovered. | terms: modeling-principle, alignment-principle, degrees-of-freedom -->

% <!-- box-family: inset -->

% <!-- inset-domain: ENGINEERING -->

% mage-box-prefix: Inset —

\begin{bookbox}{inset}{Uncertainty and unknown unknowns}
\label{inset-uncertainty}
Recall that a degree of freedom is a choice among alternatives known to be acceptable, while uncertainty concerns choices whose consequences we do not yet understand. Sometimes we have not even recognized the property or interaction that needs investigation---an unknown unknown. A model cannot directly represent a distinction we have not noticed, but organizing what we do know into entities, relationships, or states can help us discover it. We can compare alternatives, trace interactions, and examine unexpected behavior in a form simpler than the full implementation. A missing case or surprising interaction may then reveal an obligation we had not thought to specify.

This is one reason modeling can precede specification. Sometimes engineers construct a model because they know which property must hold. Sometimes they construct one to discover which properties matter.
\end{bookbox}

Modeling does not require a complete specification. Engineers can represent what they know, investigate what remains uncertain, and revise the model as they learn more. The engineering question determines what the model needs to preserve.

% <!-- point: the-question-decides-what-to-keep | Hold the system constant, change the question, and the model changes; there is no single model of the system. | terms: model-as-map, scope-of-modeling, map-and-territory -->

Ask whether two workers can process the same job at once, and the color of the web interface falls away; you need who can own work, how ownership is acquired, when it expires, and what happens after a crash. Ask instead whether one service may bypass another, and the ownership details fall away; now you need components, communication edges, and the seams a call must pass through. The system did not change. The question did. The useful reduction is not always obvious at first. Sometimes you have to look at the problem differently to see which relationships matter.

\begin{quote}
\textbf{Hold the system constant. Change the engineering question. Watch the model change.}
\end{quote}

% <!-- point: there-is-no-model-of-the-system | A model earns its place by the question it settles, not by how much of the system it represents. | terms: model-as-map, scope-of-modeling -->

A model earns its place by the question it settles, not by how much of the system it represents.

This idea also has a long history in artificial intelligence. Knowledge-representation research treats a representation as more than stored information: its form embodies commitments about what distinctions matter and affects what kinds of inference are convenient or possible\citep{davis1993kr}. MAGE applies the same principle to engineering representations.

% <!-- point: a-reduction-first-even-when-executable | Models may be encoded in types or held in an engineer’s understanding of code; modeling makes the relevant reduction available to other reasoners, and a different question may require another reduction. | terms: model-as-map, executable-source-of-truth -->

You might be wondering whether familiar programming mechanisms already do some of this modeling:

\begin{itemize}
\tightlist
\item
  \textbf{Isn't my type system a model?} My types already distinguish \texttt{UserId} from \texttt{JobId}, and the checker rejects calls that confuse them. Class definitions and interfaces also expose relationships I rely on when changing the program. Why should I make another model of information my program already contains?
\item
  \textbf{Isn't my worker pool a model?} My pool has a queue feeding eight workers, and its scheduler allows at most one job per worker. I reason about waiting jobs, occupied slots, and when another job can start. That already sounds like a model of how work gets done.
\end{itemize}

So, do these count as models? Let's look at each in turn. A model may be encoded directly in source code, or it may exist only in the engineer's understanding of that code. If an existing representation answers our question adequately, we can use it.

For the first example, type systems differ in what they can express and check. Some languages support richer relationships among types and values than others, and the checker can enforce only the distinctions those rules and declarations capture.

Even when the types contain the information we need, a model might need only a small part of the program's type structure. To understand which components depend on a central interface, we might keep that interface, a few principal classes, and the relationships among them. Fields, method bodies, and unrelated types can be left out. This view lets us examine the dependencies relevant to an interface change without searching through the whole program.

Your concept of a worker pool as a queue feeding eight workers is a mental model of the implementation. The implementation is the thing represented.\magefootnote{model-object}{Compare René Magritte's \emph{The Treachery of Images} (1929): a painting of a pipe labeled ``This is not a pipe''\citep{magritte1929}.} It may contain hundreds or thousands of lines for synchronization, error handling, cancellation, and other details. You can reason with the smaller picture because you have already reduced the implementation to a model in your mind. An agent given only the code may have to reconstruct that picture before it can answer the same question. Modeling makes your understanding explicit: it records the queue, worker slots, assignment rule, and relevant state transitions so that another human or agent can reason through them. What you have in mind when you call the worker pool a model is what you would capture.

But is that model sufficient for your engineering needs? Suppose you want to know whether and how quickly the queue will grow during a busy period. The concurrency limit cannot answer that question: two workloads can obey the same limit while one queue drains and the other grows. We need a model that relates arrival patterns, job service times, and worker capacity. It can omit the jobs' business logic and the machinery that schedules individual threads, while retaining the quantities needed to reason about backlog. The concurrency model and this queueing model describe different properties of the same running pool.

% <!-- point: explicit-models-record-accepted-intent | A model adopted as a statement of requirements records which properties later implementations must preserve; purposeful reduction makes those decisions available for human review and agent use without reproducing every code detail. | terms: modeling-principle, externalized-knowledge -->

An agent also needs to know which choices in the current implementation the engineer has accepted as requirements. A type declaration or a worker-pool implementation exposes decisions embodied in the program, but it does not automatically identify which ones must survive a change. The pool's eight-worker limit might be a resource requirement or simply today's tuning choice. An agent cannot establish its status from the number in the code. When engineers adopt a model as an authoritative statement of requirements, they record which properties later work must preserve. A descriptive model may instead record the current setting without requiring it to persist.

Intent can be recorded in source code, including type declarations or tests. We could annotate every implementation detail with its status and rationale, but modeling lets us select the properties that matter to an engineering question and put them in a form humans can understand. Engineers can review that smaller representation, and agents can use the recorded requirements when deciding which changes are acceptable. Maintaining it keeps those decisions available as the implementation changes.

\section{Constructing and Revising a Model}\label{heading:constructing-and-revising-a-model}

% <!-- section-terms: modeling-principle, degrees-of-freedom -->

These representations may be written as diagrams, tables, type declarations, or executable rules. They can support reasoning, checks, simulation, or operation. Executability changes what machinery can do with a representation; the engineering question still determines which distinctions it must preserve.\magefootnote{agent-studio-models}{Siemens\textquotesingle{} Agent Studio research illustrates this use of executable models. Agents work with formal engineering models, including SysML v2 representations, and simulation evaluates candidate electric-vehicle architectures using Simcenter Amesim. The system waits for an engineer\textquotesingle s decision at critical points\citep{liu2024agentstudio}. The 2024 research account describes potential innovations, not evidence of a deployed software factory. Its example shows a model supporting design exploration; it does not establish that a resulting implementation conforms to every represented obligation.}

% <!-- point: purposeful-reduction-is-judgment-not-algorithm | Purposeful reduction is a judgment process, not an algorithm: no procedure takes an engineering question and returns the right model, so the engineer reasons toward the boundary from more than one direction and revises it as evidence arrives. | terms: modeling-principle, degrees-of-freedom -->

Choosing those distinctions requires engineering judgment. No procedure takes an engineering question and returns the right model. The engineer reasons toward the boundary --- what the model must preserve, what it may leave free --- from more than one direction, and revises it as evidence arrives.

% <!-- point: reason-top-down-from-the-question | Reasoning top-down from the question asks what must remain true and which distinctions expressing or evaluating that obligation requires; the model preserves those, and the remaining implementation choices stay free. | terms: degrees-of-freedom, tolerance -->

One direction is top-down, from the question. What must remain true? Which distinctions are required to express or evaluate that obligation? The model preserves those; the remaining implementation choices stay free.

% <!-- point: reason-bottom-up-from-the-proposed-model | Reasoning bottom-up from the proposed model asks, for each represented distinction, what engineering question or obligation requires it; a distinction no consequential purpose justifies carrying is removed. | terms: degrees-of-freedom -->

The complementary direction is bottom-up, from the proposed model. For each distinction represented, ask what engineering question or obligation requires it. If no consequential purpose justifies carrying it, remove it. This is the discipline of parsimony, applied as a working test.

% <!-- point: revise-the-model-when-evidence-changes-the-judgment | Neither direction guarantees the boundary is right; analysis, implementation, measurement, or failure may reveal that an omitted distinction mattered or an included one did not, and the model is revised when the evidence changes the engineering judgment. | terms: degrees-of-freedom, model-as-map -->

Neither direction guarantees that the model preserves exactly what is needed. Analysis, implementation, measurement, or a failure may show that an omission was acceptable or that the model needs to change. The later examples show both possibilities.

A model may be constructed from requirements, recovered from an existing system, or developed through agent exploration. These are alternative routes rather than prescribed stages; \autoref{multiple-routes-to-engineering-models} shows their possible uses.

\begin{figure}[htbp]
\centering
\includegraphics[width=\linewidth]{book/latex/assets/multiple-routes-to-engineering-models-abe57356fe.pdf}
\caption{\emph{Multiple routes to engineering models.} A model may be constructed from requirements and intent, recovered from existing systems and experience, or discovered through agent exploration. However it arrives, the model carries engineering knowledge and the selected obligations drawn from it, and it can support three uses: engineers reason and design through it, agents are directed by it, and results are evaluated and checked against it. A model recovered from an implementation describes what exists; engineers judge whether its design is desirable and which discovered properties should become obligations. Further engineering discoveries may motivate new or revised models.}
\label{multiple-routes-to-engineering-models}
\end{figure}

% <!-- point: the-most-important-evidence-comes-from-realization | Implementing and trying a design can reveal distinctions that were not evident in the initial model. | terms: modeling-principle -->

\subsection{Learning from implementation}\label{heading:implementation-can-induce-a-model}

% <!-- point: the-direction-also-runs-in-reverse | An exploratory implementation can expose constraints, tradeoffs, state distinctions, dependencies, and bounds that were not evident in the initial design. | terms: modeling-principle, degrees-of-freedom -->

We can also learn what to model by building and trying an implementation. A prototype may expose a constraint, an unexpected interaction, or a resource limit that was not evident in the initial design. This is one benefit of commodity intelligence: engineers can investigate alternatives before deciding which relationships and requirements to preserve.

% <!-- point: a-prototype-preserves-a-realization-not-the-knowledge | A prototype retains an implementation but may not explain what engineers learned from it; record the useful relationships explicitly so later work can use the lesson without repeating the experiment. | terms: modeling-principle, externalized-knowledge -->

The prototype alone may not explain what we learned. It records one implementation, including incidental choices alongside those that mattered to the experiment. After investigating it, record the useful relationships explicitly: which constraint mattered, what bound was measured, or why one alternative worked better. That record lets later work use the lesson without repeating the experiment or reconstructing its purpose from the prototype's code.

\subsection{Judging what an agent discovers}\label{heading:judging-what-an-agent-discovers}

% <!-- point: agents-import-patterns-but-design-is-contextual | An agent can introduce useful design choices through exploratory implementation, but the engineer must judge whether they suit the present system. | terms: modeling-principle, degrees-of-freedom -->
% <!-- point: human-authority-is-not-human-invention | The agent may supply a design insight; the engineer takes responsibility for adopting it by evaluating it against the system's requirements. | terms: modeling-principle -->

An exploratory implementation can introduce useful design choices that the engineer did not specify in advance. An agent may bring a decomposition, boundary, or state distinction learned from other systems into the work. But design is contextual: a pattern that worked elsewhere may violate an obligation or erase a distinction that matters here. The engineer must evaluate the proposed design against the present system\textquotesingle s requirements. The agent may have supplied the insight, but the engineer takes responsibility for deciding whether to adopt it and which aspects later implementations should preserve.

% <!-- point: model-extraction-turns-transfer-into-local-law | Engineers record accepted design choices so later implementations can preserve the relevant relationships without reproducing incidental details. | terms: modeling-principle, governance-conversion -->
% <!-- point: not-every-extraction-becomes-law | Models can preserve both accepted obligations and descriptive observations; the representation should distinguish what later implementations must preserve from what engineers learned about the current one. | terms: modeling-principle, externalized-knowledge -->

Adopting a prospective design choice means deciding what should survive a change of implementation. Suppose an agent routes database access through a service that the user interface calls. After evaluating the design, the engineer decides that the user-interface layer should never call the database directly. A model of allowed calls between layers records that decision while leaving the internal implementation of each layer open. The same prototype may reveal that loading a particular screen currently requires three database queries. That observation is useful for investigating performance, but it need not become a requirement. A later implementation that needs fewer queries can improve performance while preserving the boundary. Modeling can preserve both kinds of knowledge, provided the representation distinguishes an accepted obligation from an observation about the current system.

% <!-- point: inference-is-not-institutional-memory | Tests can explicitly preserve accepted engineering decisions; engineers choose what the tests should require, make their purpose clear, and decide how the environment should act on failures. | terms: modeling-principle, externalized-knowledge -->

Tests can record which choices the engineer accepted as requirements. An agent may reconstruct a plausible account of how the system works and what it is intended to do, but that reconstruction alone does not establish which choices must be preserved. Test-driven development puts desired behavior into a test before implementing it, then retains the test as the implementation evolves\citep{beck1999xp}. Google\textquotesingle s Beyoncé rule expresses a related principle: write tests for the behaviors and properties you want to preserve\citep{winters2020}. Such tests make the accepted decision available to later people and agents in an executable representation.

When we encode that decision in a test, we need to preserve the requirement while leaving incidental implementation choices open. In our example, an architectural test prohibiting direct calls from the user interface to the database would record the accepted boundary. A test expecting exactly three queries would go further: it would treat the observed query count as an obligation. If the concern is performance, an upper bound on query count or response time may better express the intent. Clear test names and assertions can communicate that intent; where its reason or scope needs more explanation, record it alongside the test. The model can be expressed in the test itself without requiring a separate document.

Tests also show how Modeling (\{\{chapter:modeling\}\}) connects to Alignment (\{\{chapter:alignment\}\}). Selecting and representing the property is a modeling decision. Running the test checks the implementation against it, and a build gate that rejects a failing change enforces the result. \magevalue{chapter:one-problem-many-models} follows the cycle of implementation, observation, Modeling, and Alignment across the system lifecycle.

\section{Structure and Semantics Make Representations Machine-Operable}\label{heading:structure-and-semantics-make-representations-machine-operable}

% <!-- section-terms: modeling-principle, semantic-commitment, system-knowledge-graph, structured -->

% <!-- point: a-model-can-preserve-distinctions-without-defining-them | A representation can preserve the right distinctions while its interpreters supply most of its meaning. | terms: semantic-commitment, model-as-map -->

A model can preserve the right distinctions without completely defining what those distinctions mean. Consider a box-and-arrow diagram drawn during a design meeting. The engineers may agree that the boxes denote services and the arrows denote calls. The diagram can be useful precisely because the people in the room supply much of its meaning. They know what counts as a service, what an arrow means, which details have been omitted, and which interpretations would be nonsensical.

% <!-- point: semantic-commitment-names-the-second-axis | Semantic commitment: how much meaning the representation itself carries explicitly. | terms: semantic-commitment, structured -->

% <!-- index-def: semantic-commitment -->

That is often enough. But representations can make progressively stronger semantic commitments: they can move meaning from the interpreter into the representation itself. Boxes can become typed entities with stable identities. Arrows can become declared relations with defined source and target types. Properties can acquire units, multiplicities, constraints, or invariants. A modeling language can define rules for composition and analysis. At the stronger end, tools can determine mechanically what a representation means well enough to query, transform, check, simulate, or generate from it. Call this axis \textbf{semantic commitment}: how much of the representation\textquotesingle s meaning is carried explicitly by the representation itself, rather than supplied by the humans or agents interpreting it.

% <!-- point: semantic-commitment-is-distinct-from-purposeful-reduction | Purposeful reduction governs what the model says; semantic commitment governs how much the model itself says about what it means. | terms: semantic-commitment, model-as-map -->

This axis is distinct from purposeful reduction. Two models may preserve exactly the same engineering distinctions while making different semantic commitments. An informal component sketch and a typed system model might contain the same four components and three dependencies. They can preserve the same engineering facts while carrying different amounts of their meaning explicitly. \textbf{Purposeful reduction governs what the model says. Semantic commitment governs how much the model itself says about what it means.}

% <!-- point: stronger-semantics-buy-capabilities-and-create-obligations | Stronger semantics buy capabilities and create obligations; they are not inherently better. | terms: semantic-commitment, degrees-of-freedom -->

Greater semantic commitment buys capabilities, but it also creates obligations. More precise representations can be harder for people to author and read. Their concepts must be defined, their types and relationships maintained, and their correspondence with the implemented system preserved. Stronger semantics are therefore not inherently better. The appropriate question is what the representation must support.

% <!-- point: the-spectrum-in-concrete-artifacts | Concrete artifacts illustrate the range; the rows are examples, not maturity levels. | terms: semantic-commitment, degrees-of-freedom -->

\autoref{semantic-commitment-spectrum} illustrates the range with concrete modeling artifacts, from an informal sketch through UML 2.5.1\citep{omg2017uml} and OpenAPI/JSON Schema to SysML v2 and its semantic foundation, KerML\citep{omg2025sysml}\citep{omg2025kerml}. The rows illustrate increasing semantic commitment, whose value depends on whether the resulting precision and machine operability justify their cost. Semantic commitment and implementation freedom are separate axes. A SysML model may specify fewer implementation details than an elaborate natural-language specification while defining its chosen distinctions more precisely.

{\footnotesize
\begin{longtable}[]{@{}
  >{\raggedright\arraybackslash}p{(\linewidth - 8\tabcolsep) * \real{0.2}}
  >{\raggedright\arraybackslash}p{(\linewidth - 8\tabcolsep) * \real{0.2}}
  >{\raggedright\arraybackslash}p{(\linewidth - 8\tabcolsep) * \real{0.2}}
  >{\raggedright\arraybackslash}p{(\linewidth - 8\tabcolsep) * \real{0.2}}
  >{\raggedright\arraybackslash}p{(\linewidth - 8\tabcolsep) * \real{0.2}}@{}}
\caption[Degrees of semantic commitment, in concrete artifacts]{\emph{Semantic commitment is not model completeness.} The rows are examples, not maturity levels: a representation can define precisely what its modeled distinctions mean while deliberately leaving most implementation choices open.}\label{semantic-commitment-spectrum}\tabularnewline
\toprule\noalign{}
\begin{minipage}[b]{\linewidth}\raggedright
Semantic commitment (example)
\end{minipage} & \begin{minipage}[b]{\linewidth}\raggedright
Semantics supplied by the representation
\end{minipage} & \begin{minipage}[b]{\linewidth}\raggedright
Interpretation burden
\end{minipage} & \begin{minipage}[b]{\linewidth}\raggedright
Remaining implementation freedom
\end{minipage} & \begin{minipage}[b]{\linewidth}\raggedright
Agentic implication
\end{minipage} \\
\midrule\noalign{}
\endfirsthead
\toprule\noalign{}
\begin{minipage}[b]{\linewidth}\raggedright
Semantic commitment (example)
\end{minipage} & \begin{minipage}[b]{\linewidth}\raggedright
Semantics supplied by the representation
\end{minipage} & \begin{minipage}[b]{\linewidth}\raggedright
Interpretation burden
\end{minipage} & \begin{minipage}[b]{\linewidth}\raggedright
Remaining implementation freedom
\end{minipage} & \begin{minipage}[b]{\linewidth}\raggedright
Agentic implication
\end{minipage} \\
\midrule\noalign{}
\endhead
\bottomrule\noalign{}
\endlastfoot
\textbf{Mostly implicit} (Informal box-and-arrow architecture sketch) & Almost none beyond marks and labels; the team supplies what boxes and arrows mean & Low for readers who share the context; high for anyone who does not & Very large & Agent must infer the intended semantics and choose an implementation \\
\textbf{Notation with standardized semantics} (UML 2.5.1 state-machine diagrams) & Standardized model elements and relationships give diagrams meanings beyond their visual appearance & Requires notation knowledge and project context & Large; UML generally describes selected aspects rather than determining a program & Agent has substantially more structure to interpret, but considerable design and implementation remain \\
\textbf{Machine-readable structured model} (OpenAPI 3.1 / JSON Schema; a typed repository model such as DocAble\textquotesingle s registries and schemas) & Named entities, types, fields, constraints, references, and relationships are directly machine-readable & Structure is explicit, but domain intent remains contextual & Often still large: a schema can constrain externally relevant properties while leaving implementation open & Agent can query and transform explicit structure rather than reconstructing it from prose or code \\
\textbf{Semantically defined engineering language} (SysML v2 + KerML 1.0) & Abstract syntax, types, relationships, expressions, constraints, and semantic foundations are defined as part of the language; machine-readable interchange is standardized & Substantial language and tool knowledge, but less semantics to reconstruct & Potentially very large: precise semantics do not imply complete specification & Tools can reliably interpret and compose more of the model; agents can reason over model elements and tool-produced results rather than rediscovering their meaning \\
\end{longtable}
}

% <!-- point: the-columns-are-not-one-progression | Stronger commitment aids mechanical interpretation; it moves interpretation burden rather than removing it. | terms: semantic-commitment -->

Greater semantic commitment changes where interpretive work is done and by whom, usually making more mechanical interpretation possible. An informal sketch asks almost nothing of a reader who already shares its context and a great deal of one who does not. A semantically defined language requires language and tool knowledge up front, but leaves less meaning to be reconstructed each time the model is read.

% <!-- point: semantic-precision-is-not-generative-completeness | Semantic precision and generative completeness are different properties. | terms: semantic-commitment, degrees-of-freedom -->

This distinction matters particularly for generative implementation. Semantic precision and generative completeness are different properties. A model can have precisely defined meaning while remaining radically incomplete as a program. Conversely, a detailed prompt can say a great deal about an intended implementation while leaving much of its meaning informal. The engineering question is which decisions the engineer has made explicit and which decisions the implementer remains authorized to make.

\subsection{Giving a Model Defined Meaning}\label{heading:giving-a-model-defined-meaning}

% <!-- section-terms: semantic-commitment, model-as-map, modeling-principle -->

% <!-- point: defined-relations-make-architectural-questions-precise | Giving Calls a declared direct-invocation interpretation separates syntax from meaning; the model's purpose and scope determine whether recorded and omitted pairs describe observations or specify permitted behavior. | terms: semantic-commitment, model-as-map -->

Consider three boxes labeled Reporting, BillingAPI, and BillingDB, representing a reporting service, the billing API, and the billing database. Arrows connect Reporting to BillingAPI and BillingAPI to BillingDB. An engineer may read this as a familiar design: the reporting service obtains billing data through the API. This boundary lets the billing service control access to its data and change its schema without exposing it to every client. Yet an arrow could record a call that was observed, specify a permitted call, or express another dependency. Typing it as a relation between components tells a tool which entities may occupy its endpoints; it still leaves the relation's meaning to be defined.

Let \(C\) be the components within a declared architectural scope, including Reporting, BillingAPI, and BillingDB, and let \(\operatorname{Calls}\subseteq C\times C\). For this example, define \(\operatorname{Calls}(a,b)\) to mean that the intended architecture permits component \(a\) to invoke component \(b\) directly. Treat this relation as the complete set of permitted direct calls within that scope: a pair omitted from it is excluded by this specification. A graph of observed calls would have a different interpretation: its pairs record invocations found in a particular analysis or observation period. A missing observation establishes no prohibition and may reflect incomplete coverage. The same pair of typed component identifiers can therefore carry different claims. The engineer must choose which meaning the representation carries before using it to draw conclusions.

Suppose the obligation is that Reporting must not call BillingDB directly. We can state it over the permitted-call relation:
\[
\neg\operatorname{Calls}(\mathrm{Reporting},\mathrm{BillingDB}).
\]
A stronger obligation is that every route from Reporting to BillingDB passes through BillingAPI. Define a path as a sequence of components whose successive pairs belong to \(\operatorname{Calls}\). The requirement is then
\[
\forall p\in\operatorname{Paths}_{\operatorname{Calls}}(\mathrm{Reporting},\mathrm{BillingDB}),
\qquad \mathrm{BillingAPI}\in p.
\]
The distinction matters as the architecture grows. A route Reporting~\(\to\)~Helper~\(\to\)~BillingDB contains no direct Reporting-to-BillingDB call, but bypasses BillingAPI. Removing BillingAPI from the graph and asking whether BillingDB remains reachable from Reporting evaluates this path property. The definitions make a precise question answerable over the model while leaving algorithms and internal data structures open.

% <!-- point: transition-semantics-can-construct-intended-behavior | A partial transition function defines the intended effects of start, finish, and fail over the same four-state job lifecycle used later; undefined pairs admit no lifecycle transition. | terms: semantic-commitment, modeling-principle -->

Defined meaning can also specify how a system should operate. Consider the job lifecycle developed in \autoref{canonical-behavioral}. Its states are \(S=\{\mathrm{IDLE},\mathrm{RUNNING},\mathrm{DONE},\mathrm{FAILED}\}\), with an initial state of IDLE. Let the events be \(E=\{\mathrm{start},\mathrm{finish},\mathrm{fail}\}\). Interpret a state as the job's lifecycle status and define the effect of an event through the partial transition function \(\delta:S\times E\rightharpoonup S\):
\[
\begin{aligned}
\delta(\mathrm{IDLE},\mathrm{start}) &= \mathrm{RUNNING},\\
\delta(\mathrm{RUNNING},\mathrm{finish}) &= \mathrm{DONE},\\
\delta(\mathrm{RUNNING},\mathrm{fail}) &= \mathrm{FAILED}.
\end{aligned}
\]
These are the only defined transitions in this small specification. For every other state--event pair, \(\delta\) is undefined: the specification admits no lifecycle transition for that pair. For example, \(\delta(\mathrm{IDLE},\mathrm{finish})\) is undefined, so finish cannot take an idle job straight to completion. DONE and FAILED have no outgoing transitions in this lifecycle. The labels now denote particular statuses and events, and the arrows define their effects.

The transition function supplies behavior to implement. A job starts in IDLE, start moves it to RUNNING, and finish or fail determines its terminal outcome. An agent can implement that protocol using a database, queue, or in-memory scheduler. Alternatively, a state-machine interpreter can consult \(\delta\) to determine the next state for a given event, so the representation participates directly in execution. A test generator can use the same definition to derive event sequences and expected states: start followed by finish ends in DONE; start followed by fail ends in FAILED. These are different uses of defined meaning, with different implementation choices still open.

Both examples connect a notation to an interpretation and then to an engineering use. The structural relation defines permitted communication; the transition function defines the effects of lifecycle events. Established languages such as UML and SysML/KerML supply semantic foundations for many such constructs\citep{omg2017uml,omg2025sysml,omg2025kerml}; a small project-specific representation can also define the meaning needed for its task. Precise meaning does not, by itself, establish that an implementation conforms. The behavioral discussion below develops further analyses over executions, and \magevalue{chapter:alignment} examines how evidence and mechanisms can make conformance a condition of acceptance.

% <!-- box-family: inset -->

% <!-- inset-domain: AGENT HARNESS -->

% <!-- box-weight: aside -->

% mage-box-prefix: Inset —

\begin{bookbox}{inset}{Tools}
\label{inset-tools}
In the MAGE methodology, engineers move suitable work out of probabilistic reasoning and into engineered mechanisms. Today\textquotesingle s agent harnesses support this partly through tools. A tool gives the agent an engineered capability it can invoke: querying repository state, parsing a file format, running a compiler, inspecting a dependency graph, applying a transformation, or executing some other operation whose mechanics the agent need not reconstruct each time.

Tools change the division of labor inside the engineered system. Suppose an agent needs to determine which component owns a file. It might infer ownership from directory names, nearby code, and documentation. Or the environment might maintain an ownership model and expose a query over it. In the second case, the query supplies the recorded owner. The agent still has to decide why ownership matters and what to do with the answer, but it no longer has to infer ownership from indirect evidence. The same move applies to parsing, dependency analysis, schema validation, code transformation, and many other operations: suitable work can move from probabilistic reconstruction into mechanisms whose behavior is engineered directly. Merely exposing a tool does not guarantee that the agent will invoke it when appropriate, so consequential uses may need to be built into the surrounding workflow or checked at an enforcement boundary.

The engineering concern is therefore broader than tool use. A tool is one interface through which an engineered environment can expose knowledge or capability. The deeper design question is which parts of the work should still require the agent\textquotesingle s reasoning. As engineers record knowledge in models and provide suitable operations through queries, transformations, validators, and other mechanisms, the agent can concentrate its reasoning on the decisions that still require it.
\end{bookbox}

% <!-- point: structured-is-the-word-doing-the-work | Structure — an explicit declared shape a machine can read — is what lets a representation be traversed, checked, and joined. | terms: structured, system-knowledge-graph -->

% <!-- index-def: structured -->

Externalizing knowledge saves reconstruction only when later reasoners can recover what was externalized. Machines can do this more reliably when important entities and relations are \textbf{structured} --- given a declared shape --- rather than being recoverable only from prose. Schemas, registries, graphs, and typed records expose entities and relations that tools can traverse, validate, and join reliably.

% <!-- point: the-simplest-useful-form-is-a-graph | The simplest useful form of externalized knowledge is a graph of engineering entities and their relations; graph names the structure, not the storage. | terms: system-knowledge-graph, structured -->

% <!-- index-def: system-knowledge-graph -->

A common useful form is a graph of engineering entities and relations. Its nodes are engineering entities: services, modules, endpoints, queues, databases, requirements, tests, owners. Its edges name the relations: calls, depends on, deployed to, owned by, verified by. The graph need not live in a graph database. YAML, registries, and schemas joined on stable identifiers can represent the same structure. The word \emph{graph} names the structure, not the storage. Relevance then follows relationships: start at the entity a task names, and its neighbors come with it, so the agent reasons over a connected slice instead of a document dump.

% <!-- point: shape-does-not-distinguish-a-model-semantics-do | Shape does not distinguish representations, semantics do; semantic commitment decides how much of that meaning is declared rather than understood. | terms: system-knowledge-graph, semantic-commitment -->

Many of these representations can be interpreted as graphs of entities and relations. Their shape is not what distinguishes them; their semantics are. A dependency edge, a transition, an ownership claim, and a permission can all be drawn as arrows while answering different engineering questions. Semantic commitment determines how much of that distinction is merely understood by the reader and how much is declared in a form the engineered environment can interpret. The representation becomes useful through what those entities and relations mean, not through the fact that they happen to be drawn as nodes and edges.

% <!-- point: industry-already-lives-here | Externalizing knowledge is the lightest form of Modeling; catalogs, named obligations, reusable skills, and connected registries preserve knowledge for later work. | terms: knowledge-representation, modeling-principle, externalized-knowledge -->

Externalizing knowledge this way is the lightest form of Modeling. Spotify exposes its service catalog through tools that let agents retrieve component ownership and documentation. The catalog gives an agent an explicit place to look for that knowledge rather than leaving it to infer ownership from the implementation\citep{spotify2026constraint}. Cloudflare records standards under stable identifiers that persist as the source documents change. A particular obligation can therefore be tracked across revisions and across the systems that apply it\citep{reimann2026codex}. Shopify mines agent sessions for patterns that it incorporates into reusable skills, prompts, and defaults. Lessons from one episode can then guide later work without requiring a change to the underlying foundation model\citep{moreno2026river}. These representations differ, but each lets later agents retrieve selected project knowledge instead of reconstructing it for every task.

% <!-- point: agents-make-the-semantics-tradeoff-newly-consequential | Explicit semantics preserve meanings that later agents would otherwise have to infer without the shared context of the original team. | terms: semantic-commitment, externalized-knowledge -->

Explicit meaning helps preserve knowledge as work passes between people and agents. Newcomers still need training in institutional know-how, and some teams choose to teach that knowledge rather than record it explicitly. A team may, for example, use an informal diagram whose arrows are explained during onboarding. Agentic systems give engineers reason to revisit that choice: defining the relations in the model gives people and agents a shared reference they can consult across tasks. Engineers still decide which meanings need that precision and which can remain informal.

\section{What Structured Models Make Possible}\label{heading:representations-support-accumulating-capabilities}

% <!-- section-terms: modeling-principle, alignment-principle, capability-ladder -->

% <!-- point: capability-comes-in-rungs | Increasing semantic commitment can make additional capabilities possible — not a required sequence or maturity model. | terms: capability-ladder, structured, semantic-commitment -->

% <!-- index-def: capability-ladder -->

A model's structure and semantics determine what tools can do with it. Typed relationships support queries across related entities. Explicit properties allow checks. Other models support derivation, generation, traceability, or composition. A project may need some of these capabilities and have no use for others. Choose the representation according to the engineering work it should support (\autoref{question-to-consequence-progression}).

\begin{figure}[htbp]
\centering
\includegraphics[width=\linewidth]{book/latex/assets/question-to-consequence-progression-25598a02ec.pdf}
\caption{\emph{From engineering question to engineering consequence.} Hold the system fixed and change the question, and the useful representation changes with it. The representation determines which analyses become tractable; analysis produces evidence that can inform architecture and design, expose a property for checking, or both. Chapter 3 takes up the separate question of which obligations should be enforced.}
\label{question-to-consequence-progression}
\end{figure}

\section{The Engineering Modeling Repertoire}\label{heading:the-engineering-modeling-repertoire}

% <!-- section-terms: model-classes, projection, modeling-principle -->

% <!-- point: engineering-fields-have-accumulated-a-modeling-repertoire | Software engineers need not invent a representation from first principles each time; engineering fields have accumulated a large repertoire of models for recurring questions, each preserving different information and supporting different analyses. | terms: model-classes, modeling-principle -->

Software engineers do not need to invent a representation from first principles each time they encounter a consequential question. Engineering fields have accumulated a large repertoire of models for recurring questions: component and dependency diagrams, state machines, data-flow models, schemas, decision tables, quantitative models, provenance records, and many others. Each preserves different information and supports different analyses.

% <!-- point: read-the-repertoire-as-available-moves-not-a-taxonomy | What matters is not a representation's graphical shape but the engineering question its entities and relations preserve. | terms: model-classes, projection -->

\autoref{modeling-repertoire} shows a small portion of that repertoire. A dependency graph makes reachability and layering available for analysis. A finite-state model makes legal transitions available for state exploration or model checking. A quantitative model can expose a budget, margin, capacity, or rate to calculation. What matters is not the shape of the representation but the engineering question its entities and relations preserve.

\begin{figure}[htbp]
\centering
\includegraphics[width=\linewidth]{book/latex/assets/modeling-repertoire-f010848f36.pdf}
\caption{\emph{The engineering modeling repertoire.} Familiar representations preserve different relationships for different engineering questions. Similar graphical forms can carry different semantics: an arrow may represent a dependency, transition, flow, permission, ownership relation, or derivation.}
\label{modeling-repertoire}
\end{figure}

% <!-- point: recognize-the-question-select-the-representation | The practical skill is recognizing the engineering question and selecting a representation that preserves the properties needed to answer it; MAGE prescribes no new notation, and the inset gives the broader rule. | terms: modeling-principle, model-classes -->

The practical skill is recognizing the engineering question and selecting a representation that preserves the properties needed to answer it. MAGE prescribes no new notation. The inset gives the broader rule: reuse the modeling knowledge engineering fields have already accumulated.

% <!-- point: reuse-the-established-repertoire-choosing-is-judgment | This chapter illustrates modeling practice rather than exhausting it; once the engineering question is known, start with the representations the field has already developed — the repertoire is reusable engineering knowledge, and choosing from it remains engineering judgment. | terms: modeling-principle, model-classes -->

% <!-- box-family: inset -->

% <!-- inset-domain: ENGINEERING -->

% mage-box-prefix: Inset —

\begin{bookbox}{inset}{Reuse the engineering repertoire}
\label{inset-reuse-repertoire}
This chapter does not attempt to enumerate every useful model nor every consequential engineering question. Its purpose is to draw attention to the fact that your system has consequential questions, and that models provide a means of articulating those questions in representations that let engineers reason about those properties and give agents a succinct form in which to reason about them. The examples ahead illustrate that practice.

Start with the models, patterns, and analyses your field has already developed. In software engineering, that repertoire includes techniques from model-based systems engineering, distributed systems, and operations research. State machines capture legal transitions; dependency and data-flow graphs support reachability questions; schemas capture entities and relations; queueing and resource models expose capacity and contention; decision tables and constraint systems expose permitted combinations. Each comes with known semantics, applicable analyses, strengths, and tradeoffs.

An engineer who needs to sort data does not ordinarily invent a sorting algorithm from first principles. The engineer recognizes the problem and selects among established algorithms according to the properties and tradeoffs that matter. Modeling works the same way: recognize the engineering question, then choose a representation that preserves what is needed to answer it. Occasionally the established repertoire is inadequate and invention is required. Recognizing when that is true, and devising the new abstraction, is precisely the kind of judgment for which engineers are valuable.

Training on code, textbook-style explanations, and programming exercises can give a foundation model reusable programming knowledge\citep{chen2021codex,gunasekar2023textbooks}. An agent may already be familiar with established modeling techniques. It may recognize a state-machine pattern or suggest a queueing model without the engineer teaching the technique from first principles. The engineer still supplies the purpose of the representation, the system-specific meanings of its entities and relations, and the obligations it must preserve. Familiarity with a notation is a useful starting point; it does not establish which interpretation or assumptions fit this system.

In the MAGE methodology, engineers decide when to make that knowledge explicit in the agent's environment. They can make the relevant repertoire---or an approved subset of it---available directly to agents. DocAble's self-governance skill describes the MAGE method, a modeling repertoire, and operations agents can perform through the environment. The agent can use that guidance to recognize a situation, choose a suitable Modeling or Alignment step, and examine the resulting evidence. Selecting the appropriate technique remains an engineering judgment, but the available techniques need not be rediscovered on every task.
\end{bookbox}
